% Original course diagram. Each skip starts at the input of its own branch.
\begin{figure}[H]
\centering
\begin{tikzpicture}[
  font=\small,
  flow/.style={-{Latex[length=2mm]},thick,draw=navy},
  skip/.style={-{Latex[length=2mm]},thick,draw=coral!85!black},
  box/.style={draw=navy,fill=navy!5,rounded corners=2pt,
    minimum width=3.6cm,minimum height=0.66cm,align=center,inner sep=4pt},
  operation/.style={box,draw=teal!80!black,fill=teal!8,minimum height=0.9cm},
  add/.style={circle,draw=coral!85!black,fill=coral!10,
    minimum size=0.56cm,inner sep=0pt},
  explanation/.style={align=left,text width=4.3cm,text=navy}
]
\node[box] (input) at (0,0) {Input $X$};
\node[box] (norm1) at (0,-1.0) {LayerNorm 1};
\node[operation] (attn) at (0,-2.1)
  {Multi-head causal attention\\then dropout};
\node[add] (add1) at (0,-3.3) {$+$};
\node[box] (norm2) at (0,-4.4) {LayerNorm 2};
\node[operation,minimum width=4.3cm] (mlp) at (0,-5.5)
  {Linear $d\!\to\!4d$, GELU, linear $4d\!\to\!d$\\then dropout};
\node[add] (add2) at (0,-6.7) {$+$};
\node[box] (output) at (0,-7.65) {Output $Y$};

\draw[flow] (input) -- (norm1);
\draw[flow] (norm1) -- (attn);
\draw[flow] (attn) -- (add1);
\draw[flow] (add1) -- node[right] {$U$} (norm2);
\draw[flow] (norm2) -- (mlp);
\draw[flow] (mlp) -- (add2);
\draw[flow] (add2) -- (output);

% The first bypass carries X; the second carries the updated state U.
\draw[skip] (input.west) -- (-3,0)
  -- node[left] {$X$} (-3,-3.3) -- (add1.west);
\draw[skip] (add1.east) -- (3,-3.3)
  -- node[right] {$U$} (3,-6.7) -- (add2.east);

\node[explanation] at (6.1,-1.55)
  {Attention mixes information across visible token positions.};
\node[explanation] at (6.1,-4.7)
  {The MLP transforms each position separately, using the same weights at every
  position.};
\node[explanation] at (6.1,-7.1)
  {Both residual sums preserve the shape $B\times T\times d$.};
\end{tikzpicture}
\caption{One GPT decoder block with normalization before each branch
(pre-LayerNorm). The first residual addition produces
$U=X+\operatorname{Drop}(\operatorname{Attn}(\operatorname{LN}_1(X)))$.
The second produces
$Y=U+\operatorname{Drop}(\operatorname{MLP}(\operatorname{LN}_2(U)))$.
The two bypasses therefore carry different states, $X$ and $U$.
Dropout is active during training and disabled during evaluation.}
\label{fig:gpt-block}
\end{figure}
